\section{Measured Counter Detail}
\label{app:counters}

\paragraph{Matched profiler passes.}
All counter evidence comes from matched \texttt{msprof} op-level aggregation
passes over the same harness invocation pattern for arms 3 and 4 (same
geometry, seeds, iteration counts, and warmup). Four metric families were
collected at the large-$T$ geometry: \texttt{Memory},
\texttt{L2Cache}, \texttt{PipeUtilization}, and
\texttt{ArithmeticUtilization}, all returning successfully.

\paragraph{Kernel-level decomposition at $T{=}4096$}
(Table~\ref{tab:opdecomp}). The fused arm's \texttt{ephemeral} kernel
(single vector core) accounts for $23.8$~ms of the pass---more than all
library matmuls of either arm combined ($\approx 9.1$--$9.4$~ms)---while its
DRAM traffic is output-only. This is the direct counter evidence that (i) the
no-store mechanism holds (no core write-back appears) and (ii) the remaining
latency deficit is compute-bound single-core epilogue work, not memory
traffic---consistent with the ledger, since the two descriptor copies the
no-store arm eliminates (the $6\,\mathbf 1_{\mathrm{mat}}nr$ term of
Eq.~\eqref{eq:traffic-ledger}) are small relative to the epilogue compute
at these shapes, so removing them cannot yet move latency.

\begin{table}[h]
\centering
\footnotesize
\caption{Op-level kernel durations (ms) over one matched large-$T$
profiler pass per arm (op-level aggregation; the \texttt{Memory}-pass
durations are shown). The \texttt{ephemeral} row exists only in the fused
arm; the core round-trip rows (\texttt{Tanh}/\texttt{Muls}/\texttt{Adds}
followed by a cast to global memory) exist only in the materialized arm.}
\label{tab:opdecomp}
\begin{tabular}{lrrrr}
\toprule
Op & Arm 3 count & Arm 3 ms & Arm 4 count & Arm 4 ms \\
\midrule
Matmul (library)      & 190 & 9.01 & 143 & 9.17 \\
Ephemeral (fused gen+apply) & -- & -- & 57 & 23.82 \\
Mul (gating)          & 76  & 6.62 & 23  & 6.71 \\
Cast                  & 78  & 1.30 & 140 & 1.51 \\
SiLU                  & --  & --   & 40  & 2.75 \\
Add                   & 76  & 0.53 & 23  & 0.16 \\
Tanh (core element)   & 57  & 0.57 & --  & -- \\
Muls (core element)   & 57  & 0.42 & --  & -- \\
\bottomrule
\end{tabular}
\end{table}

\paragraph{Toolchain gaps.}
\texttt{energy\_joules} and \texttt{occupancy\_percent} are not exposed by
the \texttt{Memory}/\texttt{PipeUtilization} metric families on this CANN
build (8.5); we document this as a toolchain gap, not a skipped measurement.
Op-level L2 byte counters report zero under this \texttt{msprof} version's
summary even when L2 traffic demonstrably occurs; L2 residency at probe
geometry is therefore established through the \texttt{Memory}-pass DRAM
byte totals (both arms' working sets L2-resident; fused-arm DRAM traffic
output-only), not through L2 counters.

\paragraph{Counter-file provenance.}
Per-pass summaries were pulled from the measurement host with SHA-256
verification; the verified extracts are the evidence of record, and raw
JSON remains on the host. The publication-format record keeps
\texttt{no\_store\_claim\_ready=false} only because the energy/occupancy
fields above are unavailable from this \texttt{msprof} build---the
no-store property itself is directly counter-verified through the absence
of core write-back traffic and core round-trip ops in the fused trace.

\section{Historical Systems Tables}
\label{app:histtables}

Tables~\ref{tab:lifetimeablation}, \ref{tab:quality}, \ref{tab:scalingsweep}, and \ref{tab:statecomm} below preserve
the historical eight-process microbenchmark and the separate synthetic
all-reduce payload sizing for completeness; Section~\ref{sec:results}
explains their interpretive boundary. The conditioning ablation
(Table~\ref{tab:lifetimeablation}, main text) holds base-factor rank fixed
but skips controller/rescaling execution in the unconditioned arm, so the
arms are not compute- or state-identical; even with that caveat the
direction is unambiguous---conditioning/materialization is costly in the
unfused prototype.

\begin{table}[h]
\centering
\caption{Historical eight-process forward--backward microbenchmark at
$(d,m)=(4096,14336)$ and $(B,S)=(2,512)$.  Mean$\pm$standard deviation over
three launches; each launch ran three warmup iterations followed by ten
forward and three training timing iterations.  DDP synchronization and
optimizer steps are excluded.}
\label{tab:scalingsweep}
\begin{tabular}{rrrr}
\toprule
Rank & F+B ms & Speedup vs. dense & Tokens/s \\
\midrule
Dense & $5.528^{\dagger}$ & $1.000\times$ & 1.482M \\
64  & $3.409\pm0.100$ & $1.622\pm0.059\times$ & 2.405M \\
128 & $3.375\pm0.069$ & $1.638\pm0.030\times$ & 2.428M \\
256 & $3.429\pm0.083$ & $1.612\pm0.029\times$ & 2.390M \\
384 & $3.379\pm0.069$ & $1.636\pm0.029\times$ & 2.425M \\
\bottomrule
\end{tabular}
\\[-1mm]
{\footnotesize $^{\dagger}$Dense aggregate reconstructed from the archived
means and speedup ratios; its per-launch raw JSON and launch-level deviation
are unavailable.}
\end{table}

\begin{table}[h]
\centering
\caption{DDP-synchronized training step on 4$\times$Ascend~910B2 at the
8B-scale geometry $(d,m)=(4096,14336)$, $B{=}2,S{=}512$ (fresh ASI2 run,
2026-09-22).  Each timed step is zero-grad, forward, squared loss, backward,
\emph{with} DistributedDataParallel gradient all-reduce firing (bucket cap
25~MiB); the elementwise optimizer update is excluded from the timed region
and, when included, differs across arms (Table~\ref{tab:ddptrainopt}:
$8.97$~ms dense vs.\ $1.0$--$1.3$~ms low-rank arms).  Three repeats of
five measured iterations; barrier and device
synchronize around every iteration; max-over-ranks timing.}
\label{tab:ddptrain}
\begin{tabular}{lrr}
\toprule
Arm & Sync.\ train ms & Speedup vs.\ dense \\
\midrule
Dense & $15.230\pm0.186$ & $1.000\times$ \\
MWG r32  & $4.888\pm0.178$ & $3.116\times$ \\
MWG r64  & $4.829\pm0.199$ & $3.154\times$ \\
MWG r128 & $4.873\pm0.125$ & $3.126\times$ \\
MWG r256 & $4.803\pm0.137$ & $3.171\times$ \\
MWG r384 & $5.071\pm0.139$ & $3.003\times$ \\
MWG r512 & $5.306\pm0.095$ & $2.870\times$ \\
\bottomrule
\end{tabular}
\end{table}

\begin{table}[h]
\centering
\caption{Optimizer-inclusive DDP-synchronized training step on
4$\times$Ascend~910B2 at the 8B-scale geometry $(d,m)=(4096,14336)$,
$B{=}2,S{=}512$ (fresh ASI2 run, 2026-09-23).  Each timed step is zero-grad,
forward, squared loss, backward \emph{with} DistributedDataParallel gradient
all-reduce (bucket cap 25~MiB) \emph{and} the elementwise AdamW update
($\mathrm{lr}{=}10^{-3}$, no weight decay).  Three repeats of five measured
iterations; barrier and device synchronize around every iteration;
max-over-ranks timing.}
\label{tab:ddptrainopt}
\begin{tabular}{lrr}
\toprule
Arm & Opt-incl.\ sync.\ train ms & Speedup vs.\ dense \\
\midrule
Dense & $24.591\pm0.596$ & $1.000\times$ \\
MWG r32  & $6.095\pm0.294$ & $4.034\times$ \\
MWG r64  & $6.153\pm0.486$ & $3.997\times$ \\
MWG r96  & $6.526\pm0.938$ & $3.768\times$ \\
MWG r128 & $6.011\pm0.100$ & $4.091\times$ \\
MWG r192 & $5.940\pm0.120$ & $4.140\times$ \\
MWG r256 & $5.984\pm0.214$ & $4.109\times$ \\
MWG r384 & $5.989\pm0.201$ & $4.106\times$ \\
MWG r512 & $7.292\pm1.263$ & $3.372\times$ \\
\bottomrule
\end{tabular}
\end{table}

\paragraph{Synchronized-DDP training step.}
The historical sweeps above time local forward--backward without gradient
synchronization.  Table~\ref{tab:ddptrain} adds the synchronized
measurement: a real \texttt{DistributedDataParallel} step on four
processes,
where backward's gradient all-reduce is included in the timed region.  The
dense arm pays a $3.0\times$ synchronization premium ($15.23$ vs.\ $5.02$~ms
unsynchronized) because its $3dm$ gradient volume dominates the collective;
the low-rank arm's $r(d+m)$ gradient volume keeps the synchronized step at
$4.8$--$5.3$~ms, a $2.87\times$--$3.17\times$ speedup over dense.  The
speedup decreases with rank (r512 is lowest at $2.87\times$) as the
low-rank gradient volume approaches the dense communication cost, exactly
as the traffic model of Section~\ref{sec:analysis} predicts.

\paragraph{Optimizer-inclusive synchronized step.}
Table~\ref{tab:ddptrainopt} extends the synchronized measurement to the
complete training step, adding the elementwise AdamW update inside the timed
region.  The update does \emph{not} cancel across arms: dense pays
$8.97$~ms to update its $3dm$ parameters against $1.0$--$1.3$~ms for the
low-rank arm's $r(d{+}m)$ parameters, so the optimizer-inclusive speedup
widens to $3.37\times$--$4.14\times$ (r192 highest; r512 lowest, where its
larger gradient volume and update cost approach dense's communication
regime).

\begin{table}[h]
\centering
\caption{Element-first state accounting and separate eight-process synthetic
all-reduce latency.  Factors use fp16; the legacy controller uses fp32.  The
collective is a proxy, not an integrated DDP step or HBM measurement.}
\label{tab:statecomm}
\resizebox{\linewidth}{!}{%
\begin{tabular}{lrrrrrr}
\toprule
Model & Factor elems. (M) & Ctrl. elems. (M) & Factor MiB & Ctrl. MiB & Total MiB & AR ms \\
\midrule
Dense & 176.161 & 0 & 336.000 & 0 & 336.000 & 6.91 \\
Rank 64  & 3.539 & 1.148 & 6.750 & 4.377 & 11.127 & 1.15 \\
Rank 128 & 7.078 & 1.246 & 13.500 & 4.754 & 18.254 & 1.28 \\
Rank 256 & 14.156 & 1.444 & 27.000 & 5.508 & 32.508 & 1.42 \\
\bottomrule
\end{tabular}
}
\end{table}

\section{Eligibility Routing Details}
\label{app:routing}

The deployable per-call decision is
\begin{equation}
g_{\routebudget}(x)=
\left(\mathbf 1\{\widehat R(x)<\tau_{\routebudget}\}
+\mathbf 1\{\widehat R(x)=\tau_{\routebudget}\}h_{\routebudget}(x)\right)
\mathbf 1\{\widehat R(x)\le\Delta_{\max}\},
\label{eq:router}
\end{equation}
where the keyed tie rule $h_{\routebudget}$ depends only on a stable call
identifier, not future calls. The realized test or deployment rate need not
equal the calibration rate and must be reported with uncertainty. A scheduler
may rank only an already-present concurrent dispatch window if it needs a hard
instantaneous cap; it may not inspect future calls.

\begin{proposition}[Budgeted value and safety conditions]
\label{prop:budget}
Let $S_{\routebudget}=\{x:g_{\routebudget}(x)=1\}$ for complete dispatchable
calls. Suppose each selected generated call saves $s_r$ bytes or communication
units relative to the dense call, and let $o_r$ be the router, launch, staging,
and any dense-residency overhead paid per candidate call. The expected net
systems saving per candidate call is
\begin{equation}
\Pr[x\in S_{\routebudget}]s_r-o_r.
\end{equation}
Quality preservation additionally requires both a predeclared mean
non-inferiority margin $\delta_{\rm NI}$ and tail-risk limit:
\begin{equation}
\E[\Delta(x)\mid x\in S_{\routebudget}]\le\delta_{\rm NI},\qquad
\Pr[\Delta(x)>\epsilon\mid x\in S_{\routebudget}]\le\alpha.
\end{equation}
Mean safety alone cannot exclude severe tail failures; high quantiles and
worst-domain results must therefore accompany the aggregate.
\end{proposition}
\begin{proof}
The expected selected-call saving is
$\E[g_{\routebudget}(x)s_r]=\Pr[x\in S_{\routebudget}]s_r$; $o_r$ is paid for
every candidate call regardless of the routing decision. The two quality
inequalities are evaluation requirements rather than consequences of the
systems identity.  Their estimates require confidence intervals or calibrated
risk guarantees on disjoint data.
\end{proof}

\section{Approximation Analysis}
\label{app:approx}

\begin{proposition}[Frobenius-error FFN distortion bound]
\label{prop:approx}
View $X$ as the tokens in one group on which each approximating matrix
$\widetilde W_p$ is fixed.  Let $A=\phi(XW_{\mathrm{gate}})$,
$B=XW_{\mathrm{up}}$, $\widetilde A=\phi(X\widetilde W_{\mathrm{gate}})$, and
$\widetilde B=X\widetilde W_{\mathrm{up}}$.  Suppose
$\norm{W_p-\widetilde W_p}_F\le\delta_p$ for
$p\in\{g,u,d\}$, $\phi$ is elementwise $L_\phi$-Lipschitz,
$\norm{B}_{\max}\le M_u$, and
$\norm{\widetilde A}_{\max}\le M_g$. Then
\begin{equation}
\norm{Y-\widetilde{Y}}_F
\le \norm{W_d}_2\norm{X}_2
  \left(M_uL_\phi\delta_g+M_g\delta_u\right)
  + \norm{\widetilde A\odot\widetilde B}_2\delta_d.
\end{equation}
\end{proposition}
\begin{proof}
Write $C=A\odot B$ and
$\widetilde C=\widetilde A\odot\widetilde B$.  Then
\begin{align}
\norm{C-\widetilde C}_F
&\le \norm{(A-\widetilde A)\odot B}_F
   +\norm{\widetilde A\odot(B-\widetilde B)}_F \\
&\le M_u\norm{A-\widetilde A}_F
   +M_g\norm{B-\widetilde B}_F \\
&\le \norm{X}_2(M_uL_\phi\delta_g+M_g\delta_u),
\end{align}
using $\norm{X\Delta W}_F\le\norm{X}_2\norm{\Delta W}_F$.
Finally,
$Y-\widetilde Y=(C-\widetilde C)W_d
+\widetilde C(W_d-\widetilde W_d)$.  Apply
$\norm{PQ}_F\le\norm{P}_F\norm{Q}_2$ to the first term and
$\norm{PQ}_F\le\norm{P}_2\norm{Q}_F$ to the second.
\end{proof}

\begin{corollary}[Truncated-SVD specialization]
\label{cor:rankapprox}
For nonzero $W_p$, let
\begin{equation}
E_r(W_p)=\frac{\sum_{i=1}^{r}\sigma_i(W_p)^2}{\norm{W_p}_F^2}
\end{equation}
be the retained singular-value energy of $W_p$ at rank $r$.  For the rank-$r$
truncated SVD,
\begin{equation}
\delta_p=\norm{W_p-W_p^{(r)}}_F
=\norm{W_p}_F\sqrt{1-E_r(W_p)}.
\end{equation}
Substitution into Proposition~\ref{prop:approx} expresses each approximation-
error term through retained Frobenius energy; the activation and operator-norm
factors in the bound remain data dependent.  For a token-conditioned $q=1$
core, the proposition applies token by token and can be aggregated over groups.
By Eckart--Young--Mirsky, this SVD is
optimal among rank-at-most-$r$ approximations for unweighted fixed-matrix
Frobenius error.  A data-conditioned model may optimize activation-weighted or
task loss instead, but this elementary sensitivity bound establishes neither
learnability, generalization, router predictability, nor a systems--quality
tradeoff.
\end{corollary}



\section{Analytic Reuse Break-even Table}
\label{app:regimecall}

\begin{table}[h]
\centering
\caption{Analytic reuse break-even at $B{=}1,S{=}1024$ for the systems
geometry $(d,m,r)=(4096,14336,128)$ and the quality geometry
$(1536,8960,384)$.  $\Gamma_{\mathrm{call}}$ is the full-factor base-only
ratio of Eq.~\eqref{eq:full-factor-call};
$\rho_{\mathrm{core}}^{\mathrm{logical}}$ is the diagonal-core ratio of
Eq.~\eqref{eq:core-payload}.  Bold rows are the first $q$ at which
$\Gamma_{\mathrm{call}}<1$.  All values are optimistic logical payloads
that omit controller, source-state, and bank traffic; none is an HBM
measurement.}
\label{tab:regimecall}
\begin{tabular}{rrrrr}
\toprule
$q$ & $\lceil S/q\rceil$ & $\Gamma_{\mathrm{call}}$ (8B) & $\Gamma_{\mathrm{call}}$ (1.5B) & $\rho_{\mathrm{core}}^{\mathrm{logical}}$ (8B / 1.5B) \\
\midrule
1   & 1024 & 41.14 & 299.89 & 0.0424 / 0.3214 \\
8   & 128  & 5.14  & 37.49  & 0.0405 / 0.2964 \\
32  & 32   & 1.29  & 9.37   & 0.0402 / 0.2938 \\
\textbf{43}  & 24  & \textbf{0.96} & 7.03   & 0.0402 / 0.2935 \\
128 & 8    & 0.32  & 2.34   & 0.0402 / 0.2931 \\
\textbf{342} & 3   & 0.12  & \textbf{0.88} & 0.0402 / 0.2929 \\
512 & 2    & 0.08  & 0.59   & 0.0402 / 0.2929 \\
1024& 1    & 0.04  & 0.29   & 0.0402 / 0.2929 \\
\bottomrule
\end{tabular}
\end{table}

\section{Causal Replacement: Legacy Redundancy}
\label{app:causal}

\begin{proposition}[Redundancy of paired diagonal scales]
For $s,s'\in\R^r$, the legacy projection
$((xU)\odot s\odot s')V$ depends on $(s,s')$ only through
$z=s\odot s'$.  Hence it has no more representational power than
Eq.~\eqref{eq:diag-core} and is non-identifiable under
$(s,s')\mapsto(s\odot c,s'\oslash c)$ for componentwise nonzero $c$.
\end{proposition}
\begin{proof}
Associativity and commutativity of elementwise multiplication give the first
claim; the stated transformation preserves every component of $s\odot s'$.
\end{proof}

\paragraph{Scope of the claim.}
This proposition removes \emph{one} explicitly described redundancy: the
legacy paired-scale factorization.  It does not make the complete redesigned
parameterization identifiable.  When the factors are trained, any invertible
componentwise scaling $c\in\R^{r}$ with $U\mapsto U\,\mathrm{diag}(c)$ and
$V\mapsto V\,\mathrm{diag}(c)^{-1}$ preserves the projection
$y=(xU)\,\mathrm{diag}(z)\,V$, because componentwise diagonal matrices commute
with the diagonal core.  We therefore claim only the removal of the paired-scale
redundancy, not full identifiability of the factors.

The legacy \texttt{token\_expert\_residual} mode remains loadable only for
checkpoint compatibility.  Its soft mixture evaluates every expert and offers
no conditional-compute saving.  A $6r$ controller head cannot in general be
converted into an equivalent affine $3r$ head, so the redesigned mode requires
fresh training.  The current tests validate its information boundary and
algebra, not model quality, kernel fusion, branch skipping, or physical HBM
traffic.

\paragraph{Measured multi-seed causal quality.}
We trained the causal diagonal-core core and its controls on four disjoint
seeds (2026/2027/2028/2029) over frozen, hash-bound, numerically near-duplicate-free
train/dev/test splits (zero exact and near-duplicate cross-split overlap),
each on layer~16 of Qwen2.5-1.5B \cite{qwen2025qwen25} in fp32, 5000 steps,
and evaluated each once
on an untouched test split.
\paragraph{Activation extraction protocol.}
The splits derive from five public text suites (GSM8K test, MBPP test, the
tail of the cleaned Alpaca training set, WikiText-103 validation, and
HellaSwag validation; 1234 source records), shuffled once with a fixed seed
(20260823) and split 50/25/25 into 630/300/303 train/dev/test rows.  Deduplication
precedes the split: Unicode NFKC normalization, whitespace collapse, case
folding, masking of numeric substrings, 4-gram profiles with a 0.82
near-duplicate threshold (1 record removed, zero cross-split near-duplicate
links).  Split SHA-256 hashes and per-suite row counts are frozen in
\texttt{data/quality/manifest.json}; the test policy (no test use for
training, layer/rank/checkpoint selection, or early stopping) was written
before fitting.  Hidden states are the layer-16 FFN inputs of the instruct
model on raw text rows, max 315 training batches of batch size 1 and 128
evaluation batches, teacher targets taken from the dense FFN output; the
fp32 fitting loss is relative MSE reconstruction of that output.  This is a
single-layer frozen-teacher protocol: it does not measure perplexity,
generation quality, or any end-to-end language-model behavior, and none of
its conclusions transfer to those settings (see the evidence map,
Table~\ref{tab:evidencemap}). The dense teacher control is exact
(relative-MSE $0$, cosine $1$) in every report, the reports are marked
\texttt{evidence\_mode} under the \texttt{ASI2} label, and every checkpoint
and report fingerprint is verified. Table~\ref{tab:causalquality-app}
reports the mean$\pm$std over seeds. All quality tables in this appendix are
regenerated from a single coherent run family
(\texttt{asi2\_rank512\_control\_20260922T080042Z}) by
\texttt{scripts/make\_tables.py}, which refuses to mix rows from different
runs; the run ID is printed in every caption
(see \texttt{results/REGISTRY.md}).

\paragraph{Cross-family reconciliation.}
An earlier version of Tables 10--12 mixed rows from three distinct
executions (the rank-512 control run, the shuffled-context run of the same
date, and a 2026-09-20 closing run), which is why numbers such as
persistent r384 $0.4458$ or $0.4498$ appeared in different tables: they are
values from different run families, not contradictory measurements of one
run.  The registry binds each family (run ID, seeds, code hash, record
SHA-256) and marks every cross-family combination deprecated; the retained
families are kept as independent replications, and all paper
tables now draw exclusively from the canonical family above.  The retained
replications confirm \emph{parity}, not a directional edge: each family's
paired causal$-$persistent difference is within $0.0014$ relative MSE of
zero, but the sign of that sub-margin difference is \emph{not} stable across
families (Family~1 favors the causal core by $-0.0007$ to $-0.0008$ at
r384/r401; the closing-pull family favors the persistent control by
$+0.0005$), so no directional causal-quality claim is made.

% AUTO-GENERATED by scripts/make_tables.py -- DO NOT EDIT BY HAND
% Source: asi2_rank512_control_20260922T080042Z (results/asi2_rank512_control_20260922/)
\begin{table}[h]
\centering
\footnotesize
\caption{Causal multi-seed quality of the generated diagonal core vs.\ controls
on layer~16 of Qwen2.5-1.5B (fp32). Mean$\pm$std (sample s.d.) over four
disjoint seeds (2026--2029), 5000 steps each, disjoint hash-bound splits;
relative-MSE and cosine against the dense teacher (dense: relative-MSE $0$,
cosine $1$). All rows from the single run family
\texttt{asi2\_rank512\_control\_20260922T080042Z}.}
\label{tab:causalquality-app}
\begin{tabular}{lcccc}
\toprule
Method & Rank & Params MiB & Rel.\ MSE & Cosine \\
\midrule
Static SVD & 384 & 46.1 & $0.8369\pm0.0017$ & $0.4107\pm0.0018$ \\
Persistent low rank & 384 & 46.1 & $0.4496\pm0.0014$ & $0.7424\pm0.0009$ \\
\emph{Generated diagonal core} & 384 & 48.1 & \textbf{$0.4489\pm0.0013$} & \textbf{$0.7429\pm0.0009$} \\
\midrule
Static SVD & 401 & 48.2 & $0.8260\pm0.0017$ & $0.4236\pm0.0017$ \\
Persistent low rank & 401 & 48.2 & $0.4403\pm0.0014$ & $0.7487\pm0.0009$ \\
\emph{Generated diagonal core} & 401 & 50.2 & \textbf{$0.4395\pm0.0013$} & \textbf{$0.7492\pm0.0008$} \\
\midrule
Static SVD & 512 & 61.5 & $0.7541\pm0.0017$ & $0.5003\pm0.0016$ \\
Persistent low rank & 512 & 61.5 & $0.3833\pm0.0012$ & $0.7857\pm0.0007$ \\
\emph{Generated diagonal core} & 512 & 63.5 & \textbf{$0.3820\pm0.0012$} & \textbf{$0.7865\pm0.0007$} \\
\bottomrule
\end{tabular}
\end{table}


\paragraph{Predeclared non-inferiority margin.}
Because a non-significant difference is not evidence of equivalence, all
parity claims below are evaluated against a predeclared practical
non-inferiority margin $\delta_{\rm NI}=0.005$ relative MSE
($\approx1\%$ of the persistent control's error): the 95\% confidence
interval of the paired per-seed difference must lie inside
$[-\delta_{\rm NI},+\delta_{\rm NI}]$ for a parity claim to be made, and
seeds are treated as the unit of independence (four paired observations,
$t(3)$ intervals).  The margin was fixed before the coherent-family rebuild
and applies uniformly to every comparison in this subsection.

Three findings follow, all from the single coherent family. First, the
\emph{generated} diagonal core reaches \emph{same-rank parity} with a trained
persistent low-rank factor: its relative-MSE $0.4489\pm0.0013$ and cosine
$0.7429\pm0.0009$ match the persistent r384 control ($0.4496\pm0.0014$,
$0.7424\pm0.0009$); the paired per-seed difference is $-0.00075\pm0.00038$
with 95\% CI $[-0.00136,-0.00015]$ ($t(3){=}-3.94$, n.s.\ at $n{=}4$), and
the interval lies inside the predeclared margin
$[-0.005,+0.005]$, so the parity claim is interval-based rather than
inferred from non-significance. A core generated on the fly and
never stored thus obtains the same fitted quality as a persistent weight of the
same rank, so the ephemeral no-store contract need not trade away
representation quality at matched rank. Second, the trained low-rank path is
essential: at the same rank both the generated core and the persistent control
are markedly better than a frozen static-SVD baseline ($0.449$ vs.\ $0.837$ at
r384, $0.383$ vs.\ $0.754$ at r512). Third, at the near-equal-budget pair the
persistent r401 control (48.2~MiB) is better than the generated r384 core
(48.1~MiB): $0.4403$ vs.\ $0.4489$, $+1.95\%$ relative MSE for the generated
arm. This residual gap does \emph{not} close at matched rank: the generated
core trained at the same rank ($r{=}401$) matches persistent r401 in
reconstruction ($0.4395$ vs.\ $0.4403$, paired difference
$-0.00080\pm0.00041$), and at $r{=}512$ likewise ($0.3820$ vs.\ $0.3833$),
but because the generated core carries controller and diagonal overhead it
uses more parameters at equal rank (50.2 vs.\ 48.2~MiB at r401), so the
equal-parameter comparison retains a penalty that same-rank parity does not
remove; whether an adaptive rank allocation can close it is untested. The
causal claim is a
positive equal-rank parity result across three ranks, with an honest
$+1.95\%$ penalty at the near-equal-budget (48.1 vs.\ 48.2~MiB) comparison.

\paragraph{Shuffled-context control.}
A reviewer asked whether the generator's context-conditioning actually
matters, or whether the diagonal core's benefit is purely a parameter-count
effect.  We tested this by training a \emph{shuffled-context control}:
the same \texttt{token\_diag\_core} architecture, but with the generator's
input permuted across the batch dimension so that each token sees a random
other token's hidden state.  The actual projection still uses the real $x$;
only the generator's conditioning context is shuffled.

Table~\ref{tab:shuffled} reports the 4-seed result, drawn from the same run
family as Table~\ref{tab:causalquality-app} (a previous version of this table
mixed rows from two different runs; the mixed-source numbers, including a
$0.0029$ ``causal edge,'' are deprecated and must not be cited). Within the
coherent family the context-conditioning effect is small: at r384 the paired
causal$-$shuffled difference is $-0.00021\pm0.00003$ with 95\% CI
$[-0.00026,-0.00015]$ ($t(3){=}-12.31$, causal better in all four seeds), at
r401 $-0.00020\pm0.00004$ with 95\% CI $[-0.00026,-0.00013]$
($t(3){=}-10.13$), and at r512 $+0.00002\pm0.00002$ with 95\% CI
$[-0.00001,+0.00006]$ ($t(3){=}{+}2.50$, n.s.): about $0.05\%$ of MSE
at r384/r401, an order of magnitude smaller than the deprecated mixed-source
figure. At r512 the interval straddles zero and its upper bound
($+0.00006$) is two orders of magnitude inside $\delta_{\rm NI}$, so
\emph{the causal core shows no context-conditioning advantage over the
shuffled control at r512} --- a negative result we state plainly. The
shuffled control sits between the
causal core and the persistent baseline at r384/r401 but coincides with the
causal core at r512. Most of the diagonal-core advantage comes from the
trainable per-token correction itself, not from context-dependent generation.
This is an honest negative-leaning finding: the conditional generation
mechanism is not yet the primary quality driver at this scale, and the
controller-vs-quality question remains unfavorable to the method at this
scale.

% AUTO-GENERATED by scripts/make_tables.py -- DO NOT EDIT BY HAND
% Source: asi2_rank512_control_20260922T080042Z (results/asi2_rank512_control_20260922/)
\begin{table}[h]
\centering
\footnotesize
\caption{Shuffled-context control (4 seeds, 5000 steps each), drawn from
the same run family as Table~\ref{tab:causalquality-app}
(\texttt{asi2\_rank512\_control\_20260922T080042Z}). The shuffled generator sees a random token's
context; the real projection uses the true $x$. Paired per-seed
differences: causal$-$shuffled $= -0.00021$ (t(3)$=-12.31$) at r384,
$-0.00020$ (t(3)$=-10.13$) at r401, $+0.00002$ (t(3)$=+2.50$, n.s.) at
r512: the context-conditioning effect is $\sim$0.05\% of MSE at r384/r401
and is not detectable at r512.}
\label{tab:shuffled}
\begin{tabular}{lcc}
\toprule
Method & Rel.\ MSE & Cosine \\
\midrule
\multicolumn{3}{l}{\emph{Rank 384}} \\
Persistent low rank & $0.4496\pm0.0014$ & $0.7424\pm0.0009$ \\
Shuffled diag core & $0.4491\pm0.0013$ & $0.7428\pm0.0009$ \\
\emph{Causal diag core} & \textbf{$0.4489\pm0.0013$} & \textbf{$0.7429\pm0.0009$} \\
\midrule
\multicolumn{3}{l}{\emph{Rank 401}} \\
Persistent low rank & $0.4403\pm0.0014$ & $0.7487\pm0.0009$ \\
Shuffled diag core & $0.4397\pm0.0012$ & $0.7490\pm0.0008$ \\
\emph{Causal diag core} & \textbf{$0.4395\pm0.0013$} & \textbf{$0.7492\pm0.0008$} \\
\midrule
\multicolumn{3}{l}{\emph{Rank 512}} \\
Persistent low rank & $0.3833\pm0.0012$ & $0.7857\pm0.0007$ \\
Shuffled diag core & $0.3820\pm0.0012$ & $0.7866\pm0.0007$ \\
\emph{Causal diag core} & \textbf{$0.3820\pm0.0012$} & \textbf{$0.7865\pm0.0007$} \\
\bottomrule
\end{tabular}
\end{table}


\paragraph{Rank-512 control.}
A rank-512 persistent baseline tests whether the
generated core's quality is simply a capacity effect.  We reran the identical
protocol with a third rank, 512, for all three arms in the same run family
(four seeds, same teacher, splits, and hyperparameters; \texttt{md5} of the
harness unchanged).
Table~\ref{tab:rank512control} reports the result. Two conclusions follow.
First, same-rank parity persists at the new rank: the causal diagonal core
($0.3820\pm0.001$) matches persistent low rank ($0.3833\pm0.001$), with the
paired per-seed difference $-0.00132\pm0.00004$ and 95\% CI
$[-0.00138,-0.00125]$ (paired $t(3){=}-65.9$: statistically resolvable but
inside the predeclared $\delta_{\rm NI}=0.005$ margin, causal better in all
four seeds; the interval-based statement is practical parity, not
equivalence). Second,
and decisively, a persistent factor at rank 512 beats the generated core at
rank 384 by $14.6\%$ relative MSE ($0.3833$ vs.\ $0.4489$; paired
$t(3){=}-514$, $p{<}10^{-3}$): raw rank dominates the small same-rank
generation advantage, so the diagonal core cannot substitute for capacity,
and the ephemerality value proposition must rest on bandwidth/lifetime, not
on quality. We note that this cross-rank comparison is \emph{not}
parameter-matched (61.5 vs.\ 48.1~MiB, a $27.9\%$ storage delta); the only
near-equal-budget pair is persistent r401 vs.\ causal r384 (48.2 vs.\
48.1~MiB), which favors persistent by $+1.95\%$ relative MSE (see
Table~\ref{tab:causalquality-app} narrative). The static-SVD row ($0.7541$
at r512) confirms that training, not rank alone, drives the observed
quality.

% AUTO-GENERATED by scripts/make_tables.py -- DO NOT EDIT BY HAND
% Source: asi2_rank512_control_20260922T080042Z (results/asi2_rank512_control_20260922/)
\begin{table}[h]
\centering
\footnotesize
\caption{Rank-512 control (4 seeds, identical protocol
and run family as Table~\ref{tab:shuffled}, \texttt{asi2\_rank512\_control\_20260922T080042Z}).
At matched rank the causal core is marginally better than persistent
(paired causal$-$persistent $= -0.00132$, t(3)$=-65.94$;
the difference is statistically resolvable but is well inside any
practical non-inferiority margin). A persistent r512 factor beats the
causal r384 core by 14.6\% relative MSE
(paired t(3)$=-514$), so capacity dominates generation at
this scale.}
\label{tab:rank512control}
\begin{tabular}{lcc}
\toprule
Method (r512) & Rel.\ MSE & Cosine \\
\midrule
Static SVD & $0.7541\pm0.0017$ & $0.5003\pm0.0016$ \\
Persistent low rank & $0.3833\pm0.0012$ & $0.7857\pm0.0007$ \\
Shuffled diag core & $0.3820\pm0.0012$ & $0.7866\pm0.0007$ \\
\emph{Causal diag core} & \textbf{$0.3820\pm0.0012$} & \textbf{$0.7865\pm0.0007$} \\
\midrule
\multicolumn{3}{l}{Cross-rank (not parameter-matched at equal MiB):} \\
Persistent low rank r512 & $0.3833\pm0.0012$ & $0.7857\pm0.0007$ \\
Causal diag core r384 & $0.4489\pm0.0013$ & $0.7429\pm0.0009$ \\
\bottomrule
\end{tabular}
\end{table}


\section{Traffic Ledger}
\label{app:ledger}

\begin{equation}
L_{\mathrm{sched}}=
e_{UV}\,3r(d+m)+e_G P_G+e_K P_{\mathrm{bank}}(K)
+c_xTd+3nr+6\,\mathbf 1_{\mathrm{mat}}nr,
\label{eq:traffic-ledger}
\end{equation}
where $e_{UV},e_G,e_K$ are recorded load multiplicities, $P_G$ is controller
state, $P_{\mathrm{bank}}(K)$ is any persistent basis/expert bank,
$c_x$ counts extra source-state passes, and $\mathbf 1_{\mathrm{mat}}$ marks a
global store plus reload of the $3nr$ core values. This ledger includes
scratch-level values and is still not an HBM prediction: cache hits can suppress
durable reads, while allocator, activation, and writeback traffic add terms.

\section{Complexity Details}
\label{app:complexity}

\paragraph{Element-first parameter accounting.}
The dense FFN has $3dm$ trainable elements.  The canonical diagonal-core model
has $3r(d+m)$ persistent factor elements plus $P_G$ controller elements; the
generated $3nr$ core values are runtime intermediates, not optimizer
parameters.  With gradient element width $e_g$, the unsharded all-reduce
payload is $e_g[3r(d+m)+P_G]$.  A legacy bank of $K$ persistent rank-$r_e$
residual experts adds $3Kr_e(d+m)$ elements, and soft routing does not skip
their execution.  The canonical method sets $K=0$.  State tables report these
element classes separately before applying fp16 factor and fp32 controller
widths; no mixed-dtype total is labeled simply ``fp16.''

\paragraph{Computational complexity.}
The dense gated FFN performs three large matrix products:
$\mathcal{C}_{\mathrm{dense}} = \Theta(3Tdm)$ MACs. The associative
\mwg{} path computes $XU_r$ and then $ZV_r$ for the up, gate, and down
projections: $\mathcal{C}_{\mathrm{mwg,assoc}} = \Theta(3Tr(d+m))$. For
diagonal cores, core application adds $\Theta(3Tr)$ elementwise work and the
controller is evaluated $n=B\lceil S/q\rceil$ times under batch-isolated
grouping: $\mathcal{C}_{\mathrm{diag}}=\Theta(3Tr(d+m)+3Tr+nC_G)$. Full
factor synthesis instead emits $\Theta(3nr(d+m))$ values before the same
associative projection work. A legacy soft mixture of $K$ rank-$r_e$ experts
evaluates all $K$ residual paths; mixture weights alone do not sparsify that
work. Consequently $\Gamma(r)$ is only the leading persistent-factor compute
ratio. Small GEMMs, controller launches, and poor utilization can dominate,
as suggested by the nearly rank-flat historical timings.

\paragraph{Working-set complexity.}
The dense persistent state is $\Theta(3dm)$ elements per layer, whereas the
core method stores $\Theta(3r(d+m)+P_G)$ durable elements. Its diagonal core
has only $3r$ live values per group. Full factor synthesis would instead need
$\Theta(3r(d+m))$ live descriptor values, or a smaller tile-local subset.
These are logical working sets. Only an address-range trace and device
counters can establish that a generated core remained on chip.

\section{Break-even Proposition and Proof}
\label{app:breakeven}

\begin{proposition}[Exact full-factor reuse break-even]\label{prop:breakeven}
Let $\gamma=\Gamma(r)<1$ and
$n_{\max}=\lceil1/\gamma\rceil-1$.  Set
$g_{\max}=\lfloor n_{\max}/B\rfloor$.  Under the optimistic one-copy
accounting, a batch-isolated strict reduction exists exactly when
$g_{\max}\ge1$, and then $\Gamma_{\mathrm{call}}(r,q)<1$ exactly when
\begin{equation}
q\ge q_{\min}:=
\left\lceil\frac{S}{g_{\max}}\right\rceil.
\end{equation}
If $\gamma\ge1$ or $g_{\max}<1$, no within-example reuse group gives a
strict base-payload reduction.
\end{proposition}
\begin{proof}[Proof of Proposition~\ref{prop:breakeven}]
$\Gamma_{\mathrm{call}}<1$ is equivalent to
$B\lceil S/q\rceil<1/\gamma$, hence to
$B\lceil S/q\rceil\le n_{\max}$.  This requires
$\lceil S/q\rceil\le\lfloor n_{\max}/B\rfloor=g_{\max}$ and is equivalent to
the stated integer threshold when $g_{\max}\ge1$.
\end{proof}

\section{Lifecycle Figure}
\label{app:lifecycle}

\begin{figure}[h]
\centering
\resizebox{\linewidth}{!}{%
\begin{tikzpicture}[
  font=\footnotesize,
  node distance=7mm and 9mm,
  every node/.style={align=center},
  dur/.style={draw=black!60,rounded corners=2pt,inner sep=4pt,fill=black!4},
  gen/.style={draw=blue!60!black,rounded corners=2pt,inner sep=4pt,fill=blue!8},
  ker/.style={draw=green!50!black,rounded corners=2pt,inner sep=4pt,fill=green!8},
  rel/.style={draw=red!60!black,rounded corners=2pt,inner sep=4pt,fill=red!8},
  arr/.style={-{Stealth[length=2mm]},thick,black!70},
  lab/.style={font=\scriptsize\itshape,black!70},
  box/.style={draw=black!30,rounded corners=3pt,inner sep=6pt},
]
\node[dur] (hbm) {HBM: durable\\$U_p,V_p,\theta$};
\node[dur, right=of hbm] (l2) {L2 / cache\\staging};
\node[gen, right=of l2] (gen) {Generator\\$G_\theta(c_{\mathcal G})$};
\node[ker, right=of gen] (core) {SRAM / scratch\\$M_{p,\mathcal G}$, $r{\times}r$ diag.\\$r$ scalars / projection};
\node[ker, right=of core] (fmm) {Fused FFN\\$((XU)M)V$};
\node[rel, right=of fmm] (rel) {Overwrite};

\draw[arr] (hbm) -- node[lab,above]{persistent reads} (l2);
\draw[arr] (l2) -- (gen);
\draw[arr] (l2.north) to[bend left=18] node[lab,above]{basis reads} (fmm.north);
\draw[arr] (gen) -- node[lab,above]{emit} (core);
\draw[arr] (core) -- node[lab,above]{consume} (fmm);
\draw[arr] (fmm) -- (rel);

\begin{scope}[on background layer]
\node[box,fit=(hbm)(l2)] {};
\node[box,fit=(gen)(core)(fmm)(rel)] {};
\end{scope}

\draw[arr,dashed,red!60!black] (core.south) .. controls +(0,-8mm) and +(0,-8mm) ..
  node[lab,below,midway]{forbidden core store} (hbm.south);
\draw[arr,dashed,blue!60!black] (hbm.north) .. controls +(0,14mm) and +(0,14mm) ..
  node[lab,above,midway]{training replay / rematerialization (measured at FFN level)} (gen.north);
\end{tikzpicture}
}
\caption{Target memory hierarchy and lifecycle, not a measured implementation
of this hardware hierarchy. Persistent factors and generator weights may be
read from HBM through cache; the generated core is confined to the fused
invocation.  The red dashed path is forbidden by the inference contract.  The
blue path is the separate training-rematerialization option, measured at the
FFN level (Table~\ref{tab:lifecycle}) but not yet inside the fused kernel.}
\label{fig:lifecycle}
\end{figure}

\section{Deployment Regimes}
\label{app:regimes}

\begin{table}[h]
\centering
\caption{Deployment regimes.  The canonical method uses $K=0$ residual
experts; any $K>0$ adds durable expert reads and, under soft routing, evaluates
all experts.}
\label{tab:regimes}
\begin{tabular}{>{\raggedright\arraybackslash}p{0.17\linewidth}>{\raggedright\arraybackslash}p{0.13\linewidth}>{\raggedright\arraybackslash}p{0.27\linewidth}>{\raggedright\arraybackslash}p{0.28\linewidth}}
\toprule
Workload & Call tokens $T$ & Causal reuse $q$ & Principal constraint \\
\midrule
Prefill & $BS$, $S>1$ & $q=1$ for token-local cores; larger $q$ needs a prefix-safe shared context & Full factors require large reuse; diagonal cores keep generation small \\
Cached decode & active batch $B$, $S=1$ & Usually one core per active sequence state & Launch/controller cost can dominate; report latency and continuous-batch throughput \\
Continuous batching & variable $B$ & Grouping only if semantics permit a shared core & Report regrouping, staging, and whole-call dispatch overhead \\
Training & $BS$ & Model-dependent & Recompute-in-backward measured: $24.5\times$--$71.8\times$ saved-footprint reduction at exact gradient equality (Table~\ref{tab:lifecycle}) \\
\bottomrule
\end{tabular}
\end{table}

\section{Mechanism-Level Positioning Table}
\label{app:positioning}

Related-work positioning at the mechanism level (Section~\ref{sec:related}):

\begin{table}[h]
\centering
\caption{Closest mechanism-level comparisons.  Entries state what each
cited work \emph{reports} about its own implementation; ``n/r'' (not
reported) marks a capability the cited paper does not address and is
\emph{not} evidence of absence.  ``Measured'' reflects the four-arm probe
of Table~\ref{tab:fourarmmeasured} (fused no-store execution with matched
profiler passes); ``Target'' marks the one remaining column (branch skip)
not yet realized.}
\label{tab:positioning}
\begin{tabular}{>{\raggedright\arraybackslash}p{0.12\linewidth}>{\raggedright\arraybackslash}p{0.16\linewidth}>{\raggedright\arraybackslash}p{0.18\linewidth}ccc}
\toprule
Method & Generated object & Persistent state & Fusion & Skip & HBM counters \\
\midrule
HyperGrid & Grid-wise projection controls & Hypernetwork / base model & n/r & n/r & n/r \\
Ouroboros & Diagonal LoRA modulation & Frozen SVD-LoRA bases & n/r & n/r & n/r \\
FlashMHF & Subnetwork mixture weights & FFN subnetworks & Reported & n/r & n/r \\
MoE family & Expert index / weights & Expert banks & Varies & Reported & Varies \\
This work & Diagonal low-rank core & $U_p,V_p,G_\theta$ & Measured & Target & Partial \\
\bottomrule
\end{tabular}
\end{table}



\begin{table}[h]
\centering
\caption{Historical single-process conditioning ablation, averaged over
ranks 32--512.  Retained as design exploration only: its non-causal
controller and overlapping evaluation invalidate any quality or routing
interpretation (Appendix~\ref{app:audit-provenance}).  Speedups use each
launch family's dense control; the timing
protocol matches the eight-process sweep (three warmup, ten forward, and
three training iterations per launch).}
\label{tab:lifetimeablation}
\begin{tabular}{lrr}
\toprule
Execution arm & Mean F+B ms & Mean speedup \\
\midrule
Conditioned/materialized factors & 3.28 & $1.51\times$ \\
Persistent-factor execution & 1.53 & $3.22\times$ \\
\bottomrule
\end{tabular}
\end{table}

\begin{table}[h]
\centering
\caption{Corrected historical \emph{layer-0} reconstruction diagnostic on
Qwen2.5-1.5B.  fp32 parameter sizes are shown.}
\label{tab:quality}
\begin{tabular}{lrrrr}
\toprule
Method & Rank & Rel. MSE & Cosine & Params MiB \\
\midrule
Static SVD & 64 & 0.989 & 0.133 & 7.688 \\
Persistent low rank & 64 & 0.967 & 0.226 & 7.688 \\
Rank-scale prototype & 64 & 0.967 & 0.226 & 9.096 \\
Static SVD & 128 & 0.972 & 0.188 & 15.375 \\
Persistent low rank & 128 & 0.943 & 0.257 & 15.375 \\
Rank-scale prototype & 128 & 0.943 & 0.257 & 17.066 \\
\bottomrule
\end{tabular}
\end{table}

\section{Artifact Separation}
\label{app:artifacts}

\begin{table}[h]
\centering
\caption{Target, historical benchmark, and causal prototype are distinct.}
\label{tab:artifact-separation}
\begin{tabular}{>{\raggedright\arraybackslash}p{0.22\linewidth}>{\raggedright\arraybackslash}p{0.30\linewidth}>{\raggedright\arraybackslash}p{0.36\linewidth}}
\toprule
Artifact & Weight-like object & Evidential role \\
\midrule
Ephemeral-core target & Persistent $U_p,V_p$; generated structured $M_{p,\mathcal G}$ in scratch & Proposed inference operator; fused no-store kernel measured on device (Section~\ref{sec:systems}) \\
Systems benchmark & Persistent low-rank $U_p,V_p$ and an ordinary scale head & Measures scoped forward--backward speed; cannot attribute benefit to lifetime \\
Causal software mode & Token-local diagonal cores over persistent $U_p,V_p$ & Tests prefix/batch invariance and algebra; trained multi-seed quality measured (Table~\ref{tab:causalquality-app}) \\
\bottomrule
\end{tabular}
\end{table}

\section{Conditional Memory-Bound Speedup}
\label{app:speedup}

\begin{proposition}[Conditional memory-bound speedup]
\label{prop:speedup}
Assume $L_d=C_d+\kappa\traffic_d$ and
$L_m=C_m+\kappa\traffic_m$, where $C_d$ and $C_m$ include their respective
launch, control-generation, and compute costs. Then the generated path is
faster exactly when
\begin{equation}
C_m-C_d < \kappa(\traffic_d-\traffic_m).
\end{equation}
If both non-memory terms are negligible and
$\rho_{\mathrm{call}}=\traffic_m/\traffic_d$, the limiting speedup is
$1/\rho_{\mathrm{call}}$.
\end{proposition}
\begin{proof}
Rearranging $L_m<L_d$ gives the stated inequality. When $C_d$ and $C_m$
vanish relative to their memory terms, $L_d/L_m$ tends to
$\traffic_d/\traffic_m=1/\rho_{\mathrm{call}}$.
\end{proof}

\section{Four-Arm Measurement Figures}
\label{app:fourarmfig}

\begin{figure}[h]
\centering
\includegraphics[width=\linewidth]{figs/fourarm_ladder.pdf}
\caption{Measured four-arm evidence (data from the runs behind
Table~\ref{tab:fourarmmeasured}; medians over $200$ timed iterations
$\times$ 5 repeats). (a) At the probe geometry ($T{=}128$, working set
L2-resident by the DRAM-byte criterion of Appendix~\ref{app:counters}),
the shipped fused no-store arm sits at parity with the
materialized arm (final shipped confirmation: fused $0.823$~ms vs.\
materialized $0.851$~ms, i.e.\ $3.2\%$ faster; the $\pm3\%$ band straddles
$1.0\times$ across repeated full runs)---once host work is amortized there
is no residual device penalty for generating and consuming the core
on chip. (b) At $T{=}4096$ the logical per-projection core ($2$~MiB
fp32, $r{=}128$) exceeds the
$192$~KiB per-core UB (L2 residency inferred from DRAM byte totals, not
L2 counters; UB overflow is not evidence of L2 overflow, and the exact
L2 capacity is not separately measured); the v1-kernel fused arm is
$1.93\times$ slower, and matched
profiler passes
localize the deficit to the single-core epilogue
(Figure~\ref{fig:opdecomp}), not to core write-back traffic, which is
absent in the fused trace at both geometries. (The measured v4
multi-core rewrite of Section~\ref{sec:systems} resolves this deficit:
bit-identical to v1, $3.03\times$ faster at this geometry.)}
\label{fig:fourarm}
\end{figure}

\begin{figure}[h]
\centering
\includegraphics[width=0.62\linewidth]{figs/op_decomp.pdf}
\caption{Op-level kernel durations over matched large-$T$
\texttt{msprof} passes (Appendix~\ref{app:counters},
Table~\ref{tab:opdecomp}). The fused arm's \texttt{ephemeral}
generation+apply kernel runs on a single vector core and costs more than
all library matmuls of either arm combined, while its DRAM traffic is
output-only: the no-store property is directly visible, and the scaling
deficit is compute-bound epilogue work---the target of the v4
multi-core M-distributed epilogue design.}
\label{fig:opdecomp}
\end{figure}

\section{Four-Arm Table}
\label{app:fourarms}

\begin{table}[h]
\centering
\caption{Four-arm comparison.  The decisive lifetime contrast is arm
3 versus arm 4, not dense versus low rank.}
\label{tab:four-arms}
\begin{tabular}{clll}
\toprule
Arm & Persistent projection state & Generated object & Execution \\
\midrule
1 & Dense $W_p$ & none & Dense kernels \\
2 & Low-rank $U_p,V_p$ & none & Ordinary factor kernels \\
3 & Same $U_p,V_p,\theta$ as arm 4 & $M_{p,\mathcal G}$ & Global-memory materialization \\
4 & Same $U_p,V_p,\theta$ as arm 3 & Same $M_{p,\mathcal G}$ & Fused scratch-only, no store \\
\bottomrule
\end{tabular}
\end{table}


\section{Measured Training Lifecycle}
\label{app:training-lifecycle}

The generated diagonal core is a deterministic function of its pre-FFN input,
so backward can recompute it from a replayed input instead of materializing
and saving its intermediate activations.  We measured both lifecycles on the
NPU (Ascend 910B2, single-FFN, three seeds each): \emph{materialize-and-save}
(standard autograd) versus \emph{recompute-in-backward} (gradient
checkpointing).  Table~\ref{tab:lifecycle} reports the saved-for-backward
activation footprint, the trainer-visible peak allocator, and the
gradient-equality check.

\begin{table}[h]
\centering
\footnotesize
\caption{Measured FFN-level training lifecycle on Ascend 910B2 (three seeds
per config, all matching exactly).  Column $B,S$ is the batch and sequence
length; all configs use $(2,64)$ except the $146.98$-MiB row at $(4,128)$.
Recompute-in-backward cuts the
saved-for-backward activation footprint by $24.5\times$--$71.8\times$ relative
to materialize-and-save, with exact gradient equality ($\max$ abs diff $0.0$)
in all 12 trials.  Peak allocator is comparable at these single-FFN shapes; the
recompute path trades a small recompute-FLOP overhead for the large
saved-footprint reduction.  Single-layer scoped, matching the paper's scope;
full-model peak memory and end-to-end training benefit remain open.}
\label{tab:lifecycle}
\begin{tabular}{ccccrrr}
\toprule
$d$ & $m$ & $r$ & $B,S$ & save (MiB) & recompute (MiB) & reduction \\
\midrule
1536 & 8960 & 384 & $2,64$ & $86.04$ & $1.50$ & $57.4\times$ \\
1536 & 8960 & 384 & $4,128$ & $146.98$ & $6.00$ & $24.5\times$ \\
1536 & 8960 & 512 & $2,64$ & $107.76$ & $1.50$ & $71.8\times$ \\
2048 & 11008 & 384 & $2,64$ & $107.70$ & $2.00$ & $53.9\times$ \\
\bottomrule
\end{tabular}
\end{table}

To verify the depth-scaling claim that the saved-for-backward footprint
grows with model depth, we measured recompute-in-backward across a stack of
$N\!\in\!\{1,4,8,16\}$ matched pure FFN layers at both geometries
(Table~\ref{tab:lifecycle-depth}, Ascend 910B2).  The per-layer footprint
reduction is the geometric constant $(d{+}m)/d$---$6.38\times$ at
$d{=}2048,m{=}11008$ and $6.83\times$ at $d{=}1536,m{=}8960$---and is
\emph{independent of depth}, so the recompute advantage does not shrink in
deeper models.  Because the materialize-and-save footprint grows linearly with
$N$ while recompute stays at $N{\cdot}d$, the absolute bytes saved increase
monotonically with depth (e.g.\ at $N{=}16$ a stack spares $70$~MiB of
saved-for-backward activation per pass at $d{=}1536$, and $86$~MiB at
$d{=}2048$).  Gradient equality is
exact ($0.0$) in all 8 multi-layer configurations (single seed), so the
memory saving does not perturb training dynamics at any depth.

\begin{table}[h]
\centering
\footnotesize
\caption{Measured training lifecycle across a stack of $N$ matched FFN layers
(Ascend 910B2, $B{=}2$, $S{=}64$, fp32; \textsc{v2}, 2026-09-13).
\emph{Save} is the materialize-and-save saved-for-backward activation
footprint, \emph{recompute} that under recompute-in-backward.  The per-layer
reduction $(d{+}m)/d$ is depth-independent; the absolute footprint saved
grows linearly with depth $N$.  Gradient max-abs difference $=0.0$ in all 8
configurations (single seed).}
\label{tab:lifecycle-depth}
\begin{tabular}{cccrrr}
\toprule
$d$ & $m$ & $N$ & save (MiB) & recompute (MiB) & reduction \\
\midrule
1536 & 8960 & 1 & $5.12$ & $0.75$ & $6.83\times$ \\
1536 & 8960 & 4 & $20.50$ & $3.00$ & $6.83\times$ \\
1536 & 8960 & 8 & $41.00$ & $6.00$ & $6.83\times$ \\
1536 & 8960 & 16 & $82.00$ & $12.00$ & $6.83\times$ \\
2048 & 11008 & 1 & $6.38$ & $1.00$ & $6.38\times$ \\
2048 & 11008 & 4 & $25.50$ & $4.00$ & $6.38\times$ \\
2048 & 11008 & 8 & $51.00$ & $8.00$ & $6.38\times$ \\
2048 & 11008 & 16 & $102.00$ & $16.00$ & $6.38\times$ \\
\bottomrule
\end{tabular}
\end{table}

The saved-for-backward activation budget---the quantity that scales with
depth, batch, and sequence---is reduced by $24.5\times$--$71.8\times$.
Because gradient equality is exact, the saving does not perturb learning
dynamics.  Scope: this is a saved-tensor accounting reduction, not a measured
peak-memory reduction; at these single-FFN shapes the trainer-visible peak
allocator is comparable between the two paths, and complete-Transformer peak
memory (attention saved tensors, allocator behavior) is not measured here.  This closes the
FFN-level training-lifecycle item; the fused kernel's own training path
(retained context and generator recomputation inside the fused op, plus a
synchronized optimizer step) remains scoped future work, and the strict
principal contract stays inference-only.

\section{v4 Multi-Core Epilogue Design}
\label{app:v4design}

This appendix carries the full design detail of the v4 multi-core
epilogue summarized in Section~\ref{sec:systems}; the projections below
predate the decisive measurement and are retained as written, with the
measured decisive result recorded in the closing paragraph.

\paragraph{Failure mode of the archived v3 attempt and the v4 redesign.}
The v3 N-distribution attempt and its aicore-exception-507015 failure
under sequential launches are summarized in Section~\ref{sec:systems};
the root cause was the KFC matmul workspace shared across cores without
isolation, causing write--read collisions on back-to-back calls.  The
redesigned v4 kernel addresses both the v1 bottleneck and this failure
mode through the three coordinated mechanisms summarized in
Section~\ref{sec:systems}, detailed here:

\begin{enumerate}
  \item \emph{M-only multi-core distribution.}  Each AICore handles a
    disjoint range of $M$-tiles (tokens) with \texttt{baseN} pinned to~$r$
    (no N-split), so all cores share the same $B$ (weight) and bias tensors
    in global memory.  The host computes
    $\texttt{block\_dim}=\min(\lceil T/\texttt{baseM}\rceil,\,8)$ and each
    core's tile range from \texttt{GetBlockIdx()}/\texttt{GetBlockNum()}.
    Because the $M$-dimension is token-aligned, there is no cross-core
    synchronization and no inter-core data dependency.  \texttt{baseM} is
    halved from~128 to~64 to fit the double-buffered UB budget (see below).

  \item \emph{Per-core workspace isolation.}  The host allocates
    $\texttt{block\_dim}\times\texttt{per\_core\_ws}$ bytes of workspace,
    and each core's \texttt{REGIST\_MATMUL\_OBJ} receives a unique,
    non-overlapping slice at $\texttt{workspace}+\texttt{blockIdx}\times
    \texttt{per\_core\_ws}$.  This eliminates the shared-workspace
    collision that caused the v3 fault, at the cost of ${\sim}8\times$ more
    workspace memory (still $<\!1$~MiB total at the probe geometry).

  \item \emph{Double-buffered UB pipelining.}  All three queues (matmul
    $C$ output, activation staging, final output) use
    \texttt{BUFFER\_NUM=2}, enabling cube$\to$vector overlap within each
    core: while the vector unit processes tile~$i$'s epilogue
    (\texttt{Tanh}/\texttt{Muls}/\texttt{Adds}/\texttt{Mul}/\texttt{Cast}),
    the cube unit begins tile~$i{+}1$'s matmul into the alternate buffer.
    At \texttt{baseM}$=64$, $r=128$, $H=128$, the per-tile UB footprint is
    $2\times(C\;\text{32\,KiB}+\text{act 16\,KiB}+\text{out 16\,KiB})
    +\text{matmul staging}\sim\text{32\,KiB}\approx\text{160\,KiB}$,
    within the 192~KiB UB per core on the 910B2.
\end{enumerate}

\paragraph{How the restructuring resolves the bottleneck.}
The v1 bottleneck has three independent causes, and the v4 design addresses
each in a distinct way.  \emph{(i) Serial cube$\to$vector dependency.}  In v1,
\texttt{GetTensorC} blocks until the entire $C$ tile is in UB before
\texttt{Tanh} can begin; at $T{=}4096$ this serializes $32$ tiles.  The
$\sim$2~ms per-tile figure recorded in the v1 profiler pass is the
profiler-attributed aggregate across the full call (all tiles), not the
per-tile wall time; per-tile wall time is $\sim$60--80~$\mu$s, consistent
with the measured 2.682~ms total.  The double-buffered UB pipeline
(\texttt{BUFFER\_NUM=2}) breaks this dependency: the cube unit writes
tile~$i{+}1$ into buffer~B while the vector unit consumes tile~$i$ from
buffer~A, so the two units execute concurrently after a one-tile warm-up.
The expected per-core speedup from pipelining alone is ${\sim}1.8\times$
(the vector epilogue and cube matmul have comparable latency at the
\texttt{baseM}$=64$ tile size).  \emph{(ii) Single-core vector
utilization.}  In v1, \texttt{block\_dim}$=1$ leaves 7 of 8 AICores idle
during the epilogue.  The M-only multi-core distribution assigns
$\lceil T / (\texttt{baseM}\times 8)\rceil$ tiles per core, so all 8~cores
execute their own cube+vector pipelines in parallel.  Because the
$M$-dimension is token-aligned, no inter-core synchronization barrier is
needed---each core writes its disjoint output rows directly to global
memory.  The expected parallel speedup is ${\sim}3$--$4\times$ (sub-linear
in 8$\times$ core count due to shared-$B$ L2 read contention).
\emph{(iii) No room to double-buffer at v1's tile size.}  A previous
version of this appendix stated that v1 materializes the full $T\times r$
$C$ matrix in UB before the epilogue begins; that description was wrong and
has been regenerated from the audited kernel source.  The v1 kernel tiles
$T$ at \texttt{BASE\_M}$=128$ and holds \emph{one} $128\times r$ fp32 $C$
tile ($64$~KiB at $r{=}128$) in a depth-1 buffer at a time---the full $C$
is never resident.  The genuine constraint is that at \texttt{BASE\_M}$=128$
a second $C$ buffer does not fit: two $64$~KiB buffers are $128$~KiB before
activation staging, output, and matmul scratch, exceeding the 192~KiB UB,
so v1 cannot double-buffer and the cube and vector units remain fully
serial per tile.  The v4 design halves the tile to \texttt{baseM}$=64$
($32$~KiB per $C$ buffer), which is precisely what makes double-buffering
affordable: the per-tile UB footprint is
$64\times128\times4\,\text{B}=32$~KiB for~$C$, plus 16~KiB each for the
activation staging and output buffers, totaling ${\sim}160$~KiB with
double-buffering---within the 192~KiB budget.  The combined effect of
pipelining (${\sim}1.8\times$) and multi-core parallelism
(${\sim}3$--$4\times$) yields the projected ${\sim}5$--$7\times$ end-to-end
speedup over v1.

The no-store contract is preserved: each core reads its token slice of
$X_{\mathrm{ctrl}}$ and the rank activation from global memory, computes the
matmul and epilogue entirely in UB, and writes only the final fp16 output
to global memory.  No intermediate core tensor---the $C$ matrix, the
$\tanh$-scaled intermediate, or the bias-corrected result---is ever written
to HBM.  The v4 kernel source and benchmark harness are archived in
\texttt{custom\_ops/ephemeral\_core\_apply/torch\_ext/v4\_multicore/}.

\paragraph{Measured decisive result (2026-10-08) and remaining risk.}
The design projected ${\sim}5$--$7\times$ over v1 at $T{=}4096$
(sub-linear in core count due to shared-$B$ L2 contention), which would
bring the fused arm from $1.93\times$ \emph{slower} than materialization
to ${\sim}2.5$--$3.5\times$ \emph{faster}.  The measured decisive run on
ASI2 (hardened harness: two-phase bit-stable reference, NaN attribution,
per-config retry; 30 warmup/200 iterations/5 repeats per config) returned:
v4 is algebraically identical to v1 ($\max|\Delta_{v1-v4}|=0$ at every
config; both at the fp16 rounding floor against the reference), with
speedups of $1.13\times$ ($T{=}128$), $1.38\times$ ($512$), $1.71\times$
($1024$), $2.14\times$ ($2048$), and $3.03\times$ ($4096$).  The
direction is confirmed and the ${\sim}2\times$ viability threshold is
cleared at $T\ge2048$, but the magnitude falls short of the projection,
consistent with partial L2-contention saturation from 8~cores reading the
same $B$~tensor concurrently.  Remaining risks: (a)~UB capacity at
larger~$r$ (the 160~KiB budget is tight at $r=128$ but would overflow at
$r\ge 256$ without further tiling), and (b)~the 1000-call sequential
stability check at $T{=}4096$ (the v3 failure mode) ended in an
environment-level NaN transient in its unhardened comparison path---the
hardened five-config equivalence run above was clean, so this is
recorded as an open stability question, not a kernel defect.

